% Ismael del 11 - 19

% Reembolso tarjeta
% Reembolso bizum
% Reembolso Paypal
% Definir actividad
% Definir calendario
% Modificar actividad
% Eliminar actividad
% Listar actividades
% Anular actividad



% ======Tabla Caso de uso Número: 10 =========================
\newpage
\begin{table}[H]
    \begin{tabular}{llllllll}
    \rowcolor[HTML]{E6E6E6} % Cambiado a gris claro
    \multicolumn{1}{|l|}{Caso de Uso}                            & \multicolumn{5}{l|}{{\color[HTML]{808080} \textit{Reembolso reserva}}}                                                                                                                                                                & \multicolumn{2}{l|}{\cellcolor[HTML]{E6E6E6}{\color[HTML]{808080} \textit{CU-10}}}                                                   \\ \hline
    \multicolumn{1}{|l|}{Actores}                                & \multicolumn{7}{l|}{{\color[HTML]{808080} \textit{Cliente Regsitrado(I),Banco, Sistema PayPal}}}                                                                                                                                                                                                                                                                                               \\ \hline
    \multicolumn{1}{|l|}{Tipo}                                   & \multicolumn{7}{l|}{{\color[HTML]{808080} \textit{Primario Esencial}}}                                                                                                                                                                                                                                                                                                             \\ \hline
    \multicolumn{1}{|l|}{Referencias}                            & \multicolumn{2}{l|}{}                                                                                                   & \multicolumn{5}{p{8cm}|}{{\color[HTML]{808080} \textit{CU-11 Reembolso tarjeta (hijo), CU-12 Reembolso Bizum (hijo), CU-13 Reembolso PayPal (hijo)}}}                                                                                                                                                                                                                                    \\ \hline
    \multicolumn{1}{|l|}{Precondición}                           & \multicolumn{7}{p{10cm}|}{{\color[HTML]{808080} \textit{Un cliente registrado debe haber solicitado la cancelación de una reserva.}}}                                                                                                                                                                                                                                   \\ \hline
    \multicolumn{1}{|l|}{Poscondición}                           & \multicolumn{7}{p{10cm}|}{{\color[HTML]{808080} \textit{El dinero de la reserva debe haber sido devuelto al cliente.}}}                                                                                                                                                                                                                                                       \\ \hline
    \multicolumn{1}{|l|}{Autor}                                  & \multicolumn{1}{l|}{{\color[HTML]{808080} \textit{Ismael Sallami Moreno}}}                           & \multicolumn{2}{l|}{Fecha}                                                   & \multicolumn{1}{l|}{}                          & \multicolumn{2}{l|}{Versión}                                                      & \multicolumn{1}{l|}{1.0}                                                     \\ \hline
    \end{tabular}
\end{table}

\begin{table}[H]
    \begin{tabular}{llllllll}
    \rowcolor[HTML]{E6E6E6} % Cambiado a gris claro
    \multicolumn{8}{|l|}{\cellcolor[HTML]{F2F2F2}Propósito}                                                                                                                                                                                                                                                                                                                                                                                           \\ \hline
    \multicolumn{8}{|p{15cm}|}{{\color[HTML]{808080} \textit{Permite que un cliente registrado pueda recibir el reembolso por la cancelación de una reserva.}}}                                                                                                                                                                                                             \\ \hline
    \end{tabular}
\end{table}

\begin{table}[H]
    \begin{tabular}{llllllll}
    \rowcolor[HTML]{E6E6E6} % Cambiado a gris claro
    \multicolumn{8}{|l|}{\cellcolor[HTML]{F2F2F2}Resumen}                                                                                                                                                                                                                                                                                                                                                                                             \\ \hline
    \multicolumn{8}{|p{15cm}|}{{\color[HTML]{808080} \textit{Este caso de uso permite obtener el reembolso de la cantidad pagada por una reserva, se realizará según la forma de pago realizada: tarjeta (caso de uso Reembolso tarjeta), Bizum (caso de uso Reembolso Bizum), Paypal (caso de uso Reembolso PayPal).}}} \\ \hline
    \end{tabular}
\end{table}

\begin{table}[h!]
    \centering
    \begin{tabular}{|p{0.6cm}|p{6.5cm}|p{0.6cm}|p{6.5cm}|}
    \rowcolor[HTML]{E6E6E6} % Cambiado a gris claro
    \multicolumn{4}{|l|}{Curso Normal (Básico)} \\
    \hline
    1 & {\color[HTML]{808080} \textit{Un cliente registrado pide el reembolso de una reserva.}} & & {\color[HTML]{808080} \textit{}} \\
    \hline
     & {\color[HTML]{808080} \textit{}} & 2 & {\color[HTML]{808080} \textit{El sistema busca cuál fue el método de pago empleado en dicha reserva.}} \\
    \hline
    & & 3 & {\color[HTML]{808080} \textit{El sistema procede al reembolso equivalente al método de pago de la reserva.}} \\
    \hline
    & & 4 & {\color[HTML]{808080} \textit{El sistema notifica si el reembolso ha sido realizado correctamente o no.}} \\
    \hline
    \end{tabular}
\end{table}

\begin{table}[h!]
    \centering
    \begin{tabular}{|p{0.6cm}|p{13.4cm}|}
    \rowcolor[HTML]{E6E6E6} % Cambiado a gris claro
    \multicolumn{2}{|l|}{Cursos Alternos} \\
    \hline
    3a & {\color[HTML]{808080} \textit{Si el pago fue con PayPal, incluir CU-13 (Reembolso PayPal).}} \\
    \hline
    3b & {\color[HTML]{808080} \textit{Si el pago fue con Bizum, incluir CU-12 (Reembolso Bizum).}} \\
    \hline
    3c & {\color[HTML]{808080} \textit{Si el pago fue con tarjeta, incluir CU-11 (Reembolso tarjeta).}} \\
    \hline
    4a & {\color[HTML]{808080} \textit{Si el reembolso se realizó correctamente, se envía un justificante al cliente con la cancelación.}} \\
    \hline
    4b & {\color[HTML]{808080} \textit{Si el reembolso falló, se cancela el proceso de cancelación y termina el caso de uso.}} \\
    \hline
    \end{tabular}
\end{table}

\begin{table}[H]
    \begin{tabular}{|l
    >{\columncolor[HTML]{FFFFFF}}l l
    >{\columncolor[HTML]{FFFFFF}}l 
    >{\columncolor[HTML]{FFFFFF}}l 
    >{\columncolor[HTML]{FFFFFF}}l 
    >{\columncolor[HTML]{FFFFFF}}l 
    >{\columncolor[HTML]{FFFFFF}}l |}
    \rowcolor[HTML]{E6E6E6} % Cambiado a gris claro
    \hline
    \multicolumn{8}{|l|}{\cellcolor[HTML]{F2F2F2}Otros datos}                                                                                                                                                                                                                                                                                                        \\ \hline
    \multicolumn{1}{|l|}{Frecuencia esperada} & \multicolumn{1}{p{5cm}|}{\cellcolor[HTML]{FFFFFF}{\color[HTML]{808080} \textit{}}} & \multicolumn{1}{l|}{Rendimiento} & \multicolumn{5}{p{5cm}|}{\cellcolor[HTML]{FFFFFF}{\color[HTML]{808080} \textit{}}} \\ \hline
    \multicolumn{1}{|l|}{Importancia}         & \multicolumn{1}{p{5cm}|}{\cellcolor[HTML]{FFFFFF}{\color[HTML]{808080} \textit{}}}   & \multicolumn{1}{l|}{Urgencia}    & \multicolumn{5}{p{5cm}|}{\cellcolor[HTML]{FFFFFF}{\color[HTML]{808080} \textit{}}}                          \\ \hline
    \multicolumn{1}{|l|}{Estado}              & \multicolumn{1}{p{5cm}|}{\cellcolor[HTML]{FFFFFF}{\color[HTML]{808080} \textit{}}}                   & \multicolumn{1}{l|}{Estabilidad} & \multicolumn{5}{p{5cm}|}{\cellcolor[HTML]{FFFFFF}{\color[HTML]{808080} \textit{}}}                                            \\ \hline
    \end{tabular}
\end{table}

\begin{table}[H]
    \begin{tabular}{|l|}
    \hline
    \rowcolor[HTML]{E6E6E6} 
    \multicolumn{1}{|c|}{\textbf{Comentarios}} \\ \hline
    \multicolumn{1}{|p{15cm}|}{{\color[HTML]{808080} \textit{}}} \\ \hline
    \end{tabular}
\end{table}

% ======Tabla Caso de uso Número: 11 =========================
% Please add the following required packages to your document preamble:
% \usepackage[table,xcdraw]{xcolor}
% Beamer presentation requires \usepackage{colortbl} instead of \usepackage[table,xcdraw]{xcolor}
\newpage
\begin{table}[H]
    \begin{tabular}{llllllll}
    \hline
    \multicolumn{1}{|l|}{Caso de Uso}                            & \multicolumn{5}{l|}{{\color[HTML]{808080} \textit{Reembolso de la tarjeta}}}                                                                                                                                                                & \multicolumn{2}{l|}{\cellcolor[HTML]{E6E6E6}{\color[HTML]{808080} \textit{CU-11}}}                                                   \\ \hline
    \multicolumn{1}{|l|}{Actores}                                & \multicolumn{7}{l|}{{\color[HTML]{808080} \textit{Cliente Registrado(I), Banco}}}                                                                                                                                                                                                                                                                                               \\ \hline
    \multicolumn{1}{|l|}{Tipo}                                   & \multicolumn{7}{l|}{{\color[HTML]{808080} \textit{Primario Esencial}}}                                                                                                                                                                                                                                                                                                             \\ \hline
    \multicolumn{1}{|l|}{Referencias}                            & \multicolumn{2}{l|}{}                                                                                                   & \multicolumn{5}{l|}{{\color[HTML]{808080} \textit{CU-10 Reembolso Reserva (Padre)}}}                                                                                                                                                                                                                                    \\ \hline
    \multicolumn{1}{|l|}{Precondición}                           & \multicolumn{7}{p{10cm}|}{{\color[HTML]{808080} \textit{El cliente debe de estar registrado y debe de existir una reservar pagada mediante tarjeta.}}}                                                                                                                                                                                                                                   \\ \hline
    \multicolumn{1}{|l|}{Poscondición}                           & \multicolumn{7}{l|}{{\color[HTML]{808080} \textit{El importe de la reserva es devuelto al cliente a través de su tarjeta.}}}                                                                                                                                                                                                                                                       \\ \hline
    \multicolumn{1}{|l|}{Autor}                                  & \multicolumn{1}{l|}{{\color[HTML]{808080} \textit{Ismael Sallami Moreno}}}                           & \multicolumn{2}{l|}{Fecha}                                                   & \multicolumn{1}{l|}{}                          & \multicolumn{2}{l|}{Versión}                                                      & \multicolumn{1}{l|}{1.0}                                                     \\ \hline
    \end{tabular}
    \end{table}

    \begin{table}[H]
        \begin{tabular}{llllllll}
    \hline
    \multicolumn{8}{|l|}{\cellcolor[HTML]{F2F2F2}Propósito}                                                                                                                                                                                                                                                                                                                                                                                           \\ \hline
    \multicolumn{8}{|p{15cm}|}{{\color[HTML]{808080} \textit{Permitir al sistema procesar la devolución del importe de una reserva que fue pagada mediante tarjeta bancaria, en caso de que se cancele o anule la actividad correspondiente.}}}                                                                                                                                                                                                             \\ \hline
    \end{tabular}
    \end{table}

    \begin{table}[H]
        \begin{tabular}{llllllll}
    \hline
    \multicolumn{8}{|l|}{\cellcolor[HTML]{F2F2F2}Resumen}                                                                                                                                                                                                                                                                                                                                                                                             \\ \hline
    \multicolumn{8}{|p{15cm}|}{{\color[HTML]{808080} \textit{Este caso de uso se activa cuando un cliente ha realizado una reserva y, por alguna razón válida (como anulación de la actividad o cancelación dentro del plazo), solicita el reembolso. El sistema debe identificar que el pago se hizo por tarjeta y proceder a realizar la devolución del dinero a través del mismo método, asegurando que se cumplan las condiciones para el reembolso.}}} \\ \hline
    \end{tabular}
    \end{table}


% Please add the following required packages to your document preamble:
% \usepackage[table,xcdraw]{xcolor}
% Beamer presentation requires \usepackage{colortbl} instead of \usepackage[table,xcdraw]{xcolor}
\begin{table}[H]
    \begin{tabular}{llllllll}
    \hline
    \multicolumn{8}{|l|}{\cellcolor[HTML]{F2F2F2}Curso Normal (Básico)}                                                                                                                                                                                                                                                                          \\ \hline
    \multicolumn{1}{|c|}{1}  & \multicolumn{1}{l|}{{\color[HTML]{808080} \textit{El cliente solicita el reembolso del dinero.}}} & \multicolumn{1}{l|}{{\color[HTML]{808080} }}                    & \multicolumn{5}{l|}{{\color[HTML]{808080} }}                                                                                            \\ \hline
    \multicolumn{1}{|c|}{}   & \multicolumn{1}{l|}{{\color[HTML]{808080} \textit{}}}                                                 & \multicolumn{1}{c|}{{\color[HTML]{808080} \textit{\textbf{2}}}} & \multicolumn{5}{p{6cm}|}{{\color[HTML]{808080} \textit{El sistema verifica que la reserva fue pagada con tarjeta.}}}                         \\ \hline
    \multicolumn{1}{|l|}{}   & \multicolumn{1}{l|}{{\color[HTML]{808080} }}                                                          & \multicolumn{1}{l|}{{\color[HTML]{808080} \textbf{3}}}          & \multicolumn{5}{p{6cm}|}{{\color[HTML]{808080} \textit{El sistema valida que el reembolso es aplicable según las políticas de la empresa.}}} \\ \hline
    \multicolumn{1}{|l|}{}   & \multicolumn{1}{l|}{{\color[HTML]{808080} }}                                                          & \multicolumn{1}{c|}{{\color[HTML]{808080} \textit{4}}}          & \multicolumn{5}{p{6cm}|}{{\color[HTML]{808080} \textit{El sistema procesa el reembolso a través de la pasarela de pago.}}}                   \\ \hline
    \multicolumn{1}{|l|}{}   & \multicolumn{1}{l|}{{\color[HTML]{808080} }}                                                          & \multicolumn{1}{l|}{{\color[HTML]{808080} 5}}                   & \multicolumn{5}{p{6cm}|}{{\color[HTML]{808080} El sistema notifica al cliente el resultado del reembolso.}}                                  \\ \hline
    \end{tabular}
    \end{table}

    \begin{table}[H]
        \begin{tabular}{llllllll}
    \hline
    \multicolumn{8}{|l|}{\cellcolor[HTML]{F2F2F2}Cursos Alternos}                                                                                                                                                                                                                                                                                \\ \hline
    \multicolumn{1}{|c|}{3a} & \multicolumn{7}{p{15cm}|}{{\color[HTML]{808080} Si la reserva no cumple con las condiciones de reembolso, el sistema le informa al cliente que no se puede realizar el reembolso.}}                                                                                                                                      \\ \hline
    \multicolumn{1}{c}{}     & \multicolumn{7}{l}{}                                                                                                                                                                                                                                                                                              \\                                                                                                    &                                                                 &                            &                            &                          &                          &                        
    \end{tabular}
    \end{table}



    % Please add the following required packages to your document preamble:
    % \usepackage[table,xcdraw]{xcolor}
    % Beamer presentation requires \usepackage{colortbl} instead of \usepackage[table,xcdraw]{xcolor}
    \begin{table}[H]
    \begin{tabular}{|l
    >{\columncolor[HTML]{FFFFFF}}l l
    >{\columncolor[HTML]{FFFFFF}}l 
    >{\columncolor[HTML]{FFFFFF}}l 
    >{\columncolor[HTML]{FFFFFF}}l 
    >{\columncolor[HTML]{FFFFFF}}l 
    >{\columncolor[HTML]{FFFFFF}}l |}
    \hline
    \multicolumn{8}{|l|}{\cellcolor[HTML]{F2F2F2}Otros datos}                                                                                                                                                                                                                                                                                                        \\ \hline
    \multicolumn{1}{|l|}{Frecuencia esperada} & \multicolumn{1}{p{5cm}|}{\cellcolor[HTML]{FFFFFF}{\color[HTML]{808080} }} & \multicolumn{1}{l|}{Rendimiento} & \multicolumn{5}{p{5cm}|}{\cellcolor[HTML]{FFFFFF}{\color[HTML]{808080} }} \\ \hline
    \multicolumn{1}{|l|}{Importancia}         & \multicolumn{1}{p{5cm}|}{\cellcolor[HTML]{FFFFFF}{\color[HTML]{808080} }}   & \multicolumn{1}{l|}{Urgencia}    & \multicolumn{5}{p{5cm}|}{\cellcolor[HTML]{FFFFFF}{\color[HTML]{808080} }}                          \\ \hline
    \multicolumn{1}{|l|}{Estado}              & \multicolumn{1}{p{5cm}|}{\cellcolor[HTML]{FFFFFF}{\color[HTML]{808080} }}                   & \multicolumn{1}{l|}{Estabilidad} & \multicolumn{5}{p{5cm}|}{\cellcolor[HTML]{FFFFFF}{\color[HTML]{808080} }}                                            \\ \hline
    \end{tabular}
    \end{table}

    \begin{table}[H]
        \begin{tabular}{|l|}
        \hline
        \rowcolor[HTML]{E6E6E6} 
        \multicolumn{1}{|c|}{\textbf{Comentarios}} \\ \hline
        \multicolumn{1}{|p{15cm}|}{{\color[HTML]{808080} \textit{}}} \\ \hline
        \end{tabular}
    \end{table}



% ======Tabla Caso de uso Número: 12 =========================
\newpage
\begin{table}[H]
    \begin{tabular}{llllllll}
    \hline
    \multicolumn{1}{|l|}{Caso de Uso}                            & \multicolumn{5}{l|}{{\color[HTML]{808080} \textit{Reembolso de Bizum}}}                                                                                                                                                                & \multicolumn{2}{l|}{\cellcolor[HTML]{E6E6E6}{\color[HTML]{808080} \textit{CU-12}}}                                                   \\ \hline
    \multicolumn{1}{|l|}{Actores}                                & \multicolumn{7}{l|}{{\color[HTML]{808080} \textit{Cliente Registrado(I), Banco}}}                                                                                                                                                                                                                                                                                               \\ \hline
    \multicolumn{1}{|l|}{Tipo}                                   & \multicolumn{7}{l|}{{\color[HTML]{808080} \textit{Primario Esencial}}}                                                                                                                                                                                                                                                                                                             \\ \hline
    \multicolumn{1}{|l|}{Referencias}                            & \multicolumn{2}{l|}{}                                                                                                   & \multicolumn{5}{l|}{{\color[HTML]{808080} \textit{CU-10 Reembolso Reserva (Padre)}}}                                                                                                                                                                                                                                    \\ \hline
    \multicolumn{1}{|l|}{Precondición}                           & \multicolumn{7}{p{10cm}|}{{\color[HTML]{808080} \textit{El cliente debe de estar registrado y debe de existir una reserva pagada mediante Bizum.}}}                                                                                                                                                                                                                                   \\ \hline
    \multicolumn{1}{|l|}{Poscondición}                           & \multicolumn{7}{l|}{{\color[HTML]{808080} \textit{El importe de la reserva es devuelto al cliente a través de Bizum.}}}                                                                                                                                                                                                                                                       \\ \hline
    \multicolumn{1}{|l|}{Autor}                                  & \multicolumn{1}{l|}{{\color[HTML]{808080} \textit{Ismael Sallami Moreno}}}                           & \multicolumn{2}{l|}{Fecha}                                                   & \multicolumn{1}{l|}{}                          & \multicolumn{2}{l|}{Versión}                                                      & \multicolumn{1}{l|}{1.0}                                                     \\ \hline
    \end{tabular}
    \end{table}

    \begin{table}[H]
        \begin{tabular}{llllllll}
    \hline
    \multicolumn{8}{|l|}{\cellcolor[HTML]{F2F2F2}Propósito}                                                                                                                                                                                                                                                                                                                                                                                           \\ \hline
    \multicolumn{8}{|p{15cm}|}{{\color[HTML]{808080} \textit{Permitir al sistema procesar la devolución del importe de una reserva que fue pagada mediante Bizum, en caso de que se cancele o anule la actividad correspondiente.}}}                                                                                                                                                                                                             \\ \hline
    \end{tabular}
    \end{table}

    \begin{table}[H]
        \begin{tabular}{llllllll}
    \hline
    \multicolumn{8}{|l|}{\cellcolor[HTML]{F2F2F2}Resumen}                                                                                                                                                                                                                                                                                                                                                                                             \\ \hline
    \multicolumn{8}{|p{15cm}|}{{\color[HTML]{808080} \textit{Este caso de uso se activa cuando un cliente ha realizado una reserva y, por alguna razón válida (como anulación de la actividad o cancelación dentro del plazo), solicita el reembolso. El sistema debe identificar que el pago se hizo por Bizum y proceder a realizar la devolución del dinero a través del mismo método, asegurando que se cumplan las condiciones para el reembolso.}}} \\ \hline
    \end{tabular}
    \end{table}

\begin{table}[H]
    \begin{tabular}{llllllll}
    \hline
    \multicolumn{8}{|l|}{\cellcolor[HTML]{F2F2F2}Curso Normal (Básico)}                                                                                                                                                                                                                                                                          \\ \hline
    \multicolumn{1}{|c|}{1}  & \multicolumn{1}{l|}{{\color[HTML]{808080} \textit{El cliente solicita el reembolso de una reserva.}}} & \multicolumn{1}{l|}{{\color[HTML]{808080} }}                    & \multicolumn{5}{l|}{{\color[HTML]{808080} }}                                                                                            \\ \hline
    \multicolumn{1}{|c|}{}   & \multicolumn{1}{l|}{{\color[HTML]{808080} \textit{}}}                                                 & \multicolumn{1}{c|}{{\color[HTML]{808080} \textit{\textbf{2}}}} & \multicolumn{5}{p{6cm}|}{{\color[HTML]{808080} \textit{El sistema verifica que la reserva fue pagada con Bizum.}}}                         \\ \hline
    \multicolumn{1}{|l|}{}   & \multicolumn{1}{l|}{{\color[HTML]{808080} }}                                                          & \multicolumn{1}{l|}{{\color[HTML]{808080} \textbf{3}}}          & \multicolumn{5}{p{6cm}|}{{\color[HTML]{808080} \textit{El sistema valida que el reembolso es aplicable según las políticas de la empresa.}}} \\ \hline
    \multicolumn{1}{|l|}{}   & \multicolumn{1}{l|}{{\color[HTML]{808080} }}                                                          & \multicolumn{1}{c|}{{\color[HTML]{808080} \textit{4}}}          & \multicolumn{5}{p{6cm}|}{{\color[HTML]{808080} \textit{El sistema procesa el reembolso a través de Bizum.}}}                   \\ \hline
    \multicolumn{1}{|l|}{}   & \multicolumn{1}{l|}{{\color[HTML]{808080} }}                                                          & \multicolumn{1}{l|}{{\color[HTML]{808080} 5}}                   & \multicolumn{5}{p{6cm}|}{{\color[HTML]{808080} El sistema notifica al cliente el resultado del reembolso.}}                                  \\ \hline
    \end{tabular}
    \end{table}

    \begin{table}[H]
        \begin{tabular}{llllllll}
    \hline
    \multicolumn{8}{|l|}{\cellcolor[HTML]{F2F2F2}Cursos Alternos}                                                                                                                                                                                                                                                                                \\ \hline
    \multicolumn{1}{|c|}{3a} & \multicolumn{7}{p{15cm}|}{{\color[HTML]{808080} Si la reserva no cumple con las condiciones de reembolso, el sistema le informa al cliente que no se puede realizar el reembolso.}}                                                                                                                                      \\ \hline
    \multicolumn{1}{c}{}     & \multicolumn{7}{l}{}                                                                                                                                                                                                                                                                                              \\                                                                                                    &                                                                 &                            &                            &                          &                          &                        
    \end{tabular}
    \end{table}

    % Please add the following required packages to your document preamble:
    % \usepackage[table,xcdraw]{xcolor}
    % Beamer presentation requires \usepackage{colortbl} instead of \usepackage[table,xcdraw]{xcolor}
    \begin{table}[H]
        \begin{tabular}{|l
        >{\columncolor[HTML]{FFFFFF}}l l
        >{\columncolor[HTML]{FFFFFF}}l 
        >{\columncolor[HTML]{FFFFFF}}l 
        >{\columncolor[HTML]{FFFFFF}}l 
        >{\columncolor[HTML]{FFFFFF}}l 
        >{\columncolor[HTML]{FFFFFF}}l |}
        \hline
        \multicolumn{8}{|l|}{\cellcolor[HTML]{F2F2F2}Otros datos}                                                                                                                                                                                                                                                                                                        \\ \hline
        \multicolumn{1}{|l|}{Frecuencia esperada} & \multicolumn{1}{p{5cm}|}{\cellcolor[HTML]{FFFFFF}{\color[HTML]{808080} }} & \multicolumn{1}{l|}{Rendimiento} & \multicolumn{5}{p{5cm}|}{\cellcolor[HTML]{FFFFFF}{\color[HTML]{808080} }} \\ \hline
        \multicolumn{1}{|l|}{Importancia}         & \multicolumn{1}{p{5cm}|}{\cellcolor[HTML]{FFFFFF}{\color[HTML]{808080} }}   & \multicolumn{1}{l|}{Urgencia}    & \multicolumn{5}{p{5cm}|}{\cellcolor[HTML]{FFFFFF}{\color[HTML]{808080} }}                          \\ \hline
        \multicolumn{1}{|l|}{Estado}              & \multicolumn{1}{p{5cm}|}{\cellcolor[HTML]{FFFFFF}{\color[HTML]{808080} }}                   & \multicolumn{1}{l|}{Estabilidad} & \multicolumn{5}{p{5cm}|}{\cellcolor[HTML]{FFFFFF}{\color[HTML]{808080} }}                                            \\ \hline
        \end{tabular}
        \end{table}

        \begin{table}[H]
            \begin{tabular}{|l|}
            \hline
            \rowcolor[HTML]{E6E6E6} 
            \multicolumn{1}{|c|}{\textbf{Comentarios}} \\ \hline
            \multicolumn{1}{|p{15cm}|}{{\color[HTML]{808080} \textit{}}} \\ \hline
            \end{tabular}
        \end{table}


    % ======Tabla Caso de uso Número: 13 =========================
    \newpage
    \begin{table}[H]
        \begin{tabular}{llllllll}
        \hline
        \multicolumn{1}{|l|}{Caso de Uso}                            & \multicolumn{5}{l|}{{\color[HTML]{808080} \textit{Reembolso de PayPal}}}                                                                                                                                                                & \multicolumn{2}{l|}{\cellcolor[HTML]{E6E6E6}{\color[HTML]{808080} \textit{CU-13}}}                                                   \\ \hline
        \multicolumn{1}{|l|}{Actores}                                & \multicolumn{7}{l|}{{\color[HTML]{808080} \textit{Cliente Registrado(I), Banco y Sistema Paypal}}}                                                                                                                                                                                                                                                                                               \\ \hline
        \multicolumn{1}{|l|}{Tipo}                                   & \multicolumn{7}{l|}{{\color[HTML]{808080} \textit{Primario Esencial}}}                                                                                                                                                                                                                                                                                                             \\ \hline
        \multicolumn{1}{|l|}{Referencias}                            & \multicolumn{2}{l|}{}                                                                                                   & \multicolumn{5}{l|}{{\color[HTML]{808080} \textit{CU-10 Reembolso Reserva (Padre)}}}                                                                                                                                                                                                                                    \\ \hline
        \multicolumn{1}{|l|}{Precondición}                           & \multicolumn{7}{p{10cm}|}{{\color[HTML]{808080} \textit{El cliente debe de estar registrado y debe de existir una reserva pagada mediante PayPal.}}}                                                                                                                                                                                                                                   \\ \hline
        \multicolumn{1}{|l|}{Poscondición}                           & \multicolumn{7}{l|}{{\color[HTML]{808080} \textit{El importe de la reserva es devuelto al cliente a través de PayPal.}}}                                                                                                                                                                                                                                                       \\ \hline
        \multicolumn{1}{|l|}{Autor}                                  & \multicolumn{1}{l|}{{\color[HTML]{808080} \textit{Ismael Sallami Moreno}}}                           & \multicolumn{2}{l|}{Fecha}                                                   & \multicolumn{1}{l|}{}                          & \multicolumn{2}{l|}{Versión}                                                      & \multicolumn{1}{l|}{1.0}                                                     \\ \hline
        \end{tabular}
        \end{table}

        \begin{table}[H]
            \begin{tabular}{llllllll}
        \hline
        \multicolumn{8}{|l|}{\cellcolor[HTML]{F2F2F2}Propósito}                                                                                                                                                                                                                                                                                                                                                                                           \\ \hline
        \multicolumn{8}{|p{15cm}|}{{\color[HTML]{808080} \textit{Permitir al sistema procesar la devolución del importe de una reserva que fue pagada mediante PayPal, en caso de que se cancele o anule la actividad correspondiente.}}}                                                                                                                                                                                                             \\ \hline
        \end{tabular}
        \end{table}

        \begin{table}[H]
            \begin{tabular}{llllllll}
        \hline
        \multicolumn{8}{|l|}{\cellcolor[HTML]{F2F2F2}Resumen}                                                                                                                                                                                                                                                                                                                                                                                             \\ \hline
        \multicolumn{8}{|p{15cm}|}{{\color[HTML]{808080} \textit{Este caso de uso se activa cuando un cliente ha realizado una reserva y, por alguna razón válida (como anulación de la actividad o cancelación dentro del plazo), solicita el reembolso. El sistema debe identificar que el pago se hizo por PayPal y proceder a realizar la devolución del dinero a través del mismo método, asegurando que se cumplan las condiciones para el reembolso.}}} \\ \hline
        \end{tabular}
        \end{table}

    \begin{table}[H]
        \begin{tabular}{llllllll}
        \hline
        \multicolumn{8}{|l|}{\cellcolor[HTML]{F2F2F2}Curso Normal (Básico)}                                                                                                                                                                                                                                                                          \\ \hline
        \multicolumn{1}{|c|}{1}  & \multicolumn{1}{l|}{{\color[HTML]{808080} \textit{El cliente solicita el reembolso de una reserva.}}} & \multicolumn{1}{l|}{{\color[HTML]{808080} }}                    & \multicolumn{5}{l|}{{\color[HTML]{808080} }}                                                                                            \\ \hline
        \multicolumn{1}{|c|}{}   & \multicolumn{1}{l|}{{\color[HTML]{808080} \textit{}}}                                                 & \multicolumn{1}{c|}{{\color[HTML]{808080} \textit{\textbf{2}}}} & \multicolumn{5}{p{6cm}|}{{\color[HTML]{808080} \textit{El sistema verifica que la reserva fue pagada con PayPal.}}}                         \\ \hline
        \multicolumn{1}{|l|}{}   & \multicolumn{1}{l|}{{\color[HTML]{808080} }}                                                          & \multicolumn{1}{l|}{{\color[HTML]{808080} \textbf{3}}}          & \multicolumn{5}{p{6cm}|}{{\color[HTML]{808080} \textit{El sistema valida que el reembolso es aplicable según las políticas de la empresa.}}} \\ \hline
        \multicolumn{1}{|l|}{}   & \multicolumn{1}{l|}{{\color[HTML]{808080} }}                                                          & \multicolumn{1}{c|}{{\color[HTML]{808080} \textit{4}}}          & \multicolumn{5}{p{6cm}|}{{\color[HTML]{808080} \textit{El sistema procesa el reembolso a través de PayPal.}}}                   \\ \hline
        \multicolumn{1}{|l|}{}   & \multicolumn{1}{l|}{{\color[HTML]{808080} }}                                                          & \multicolumn{1}{l|}{{\color[HTML]{808080} 5}}                   & \multicolumn{5}{p{6cm}|}{{\color[HTML]{808080} El sistema notifica al cliente el resultado del reembolso.}}                                  \\ \hline
        \end{tabular}
        \end{table}

        \begin{table}[H]
            \begin{tabular}{llllllll}
        \hline
        \multicolumn{8}{|l|}{\cellcolor[HTML]{F2F2F2}Cursos Alternos}                                                                                                                                                                                                                                                                                \\ \hline
        \multicolumn{1}{|c|}{3a} & \multicolumn{7}{p{15cm}|}{{\color[HTML]{808080} Si la reserva no cumple con las condiciones de reembolso, el sistema le informa al cliente que no se puede realizar el reembolso.}}                                                                                                                                      \\ \hline
        \multicolumn{1}{c}{}     & \multicolumn{7}{l}{}                                                                                                                                                                                                                                                                                              \\                                                                                                    &                                                                 &                            &                            &                          &                          &                        
        \end{tabular}
        \end{table}

        \begin{table}[H]
            \begin{tabular}{|l
            >{\columncolor[HTML]{FFFFFF}}l l
            >{\columncolor[HTML]{FFFFFF}}l 
            >{\columncolor[HTML]{FFFFFF}}l 
            >{\columncolor[HTML]{FFFFFF}}l 
            >{\columncolor[HTML]{FFFFFF}}l 
            >{\columncolor[HTML]{FFFFFF}}l |}
            \hline
            \multicolumn{8}{|l|}{\cellcolor[HTML]{F2F2F2}Otros datos}                                                                                                                                                                                                                                                                                                        \\ \hline
            \multicolumn{1}{|l|}{Frecuencia esperada} & \multicolumn{1}{p{5cm}|}{\cellcolor[HTML]{FFFFFF}{\color[HTML]{808080} }} & \multicolumn{1}{l|}{Rendimiento} & \multicolumn{5}{p{5cm}|}{\cellcolor[HTML]{FFFFFF}{\color[HTML]{808080} }} \\ \hline
            \multicolumn{1}{|l|}{Importancia}         & \multicolumn{1}{p{5cm}|}{\cellcolor[HTML]{FFFFFF}{\color[HTML]{808080} }}   & \multicolumn{1}{l|}{Urgencia}    & \multicolumn{5}{p{5cm}|}{\cellcolor[HTML]{FFFFFF}{\color[HTML]{808080} }}                          \\ \hline
            \multicolumn{1}{|l|}{Estado}              & \multicolumn{1}{p{5cm}|}{\cellcolor[HTML]{FFFFFF}{\color[HTML]{808080} }}                   & \multicolumn{1}{l|}{Estabilidad} & \multicolumn{5}{p{5cm}|}{\cellcolor[HTML]{FFFFFF}{\color[HTML]{808080} }}                                            \\ \hline
            \end{tabular}
            \end{table}

            \begin{table}[H]
                \begin{tabular}{|l|}
                \hline
                \rowcolor[HTML]{E6E6E6} 
                \multicolumn{1}{|c|}{\textbf{Comentarios}} \\ \hline
                \multicolumn{1}{|p{15cm}|}{{\color[HTML]{808080} \textit{}}} \\ \hline
                \end{tabular}
            \end{table}
    
    
    
    % ======Tabla Caso de uso Número: 14 =========================
    \chapter{Gestión de Actividades}
    \begin{table}[H]
        \begin{tabular}{llllllll}
        \hline
        \multicolumn{1}{|l|}{Caso de Uso}                            & \multicolumn{5}{l|}{{\color[HTML]{808080} \textit{Definir actividad}}}                                                                                                                                                                & \multicolumn{2}{l|}{\cellcolor[HTML]{E6E6E6}{\color[HTML]{808080} \textit{CU-14}}}                                                   \\ \hline
        \multicolumn{1}{|l|}{Actores}                                & \multicolumn{7}{l|}{{\color[HTML]{808080} \textit{Empleado(I)}}}                                                                                                                                                                                                                                                                                               \\ \hline
        \multicolumn{1}{|l|}{Tipo}                                   & \multicolumn{7}{l|}{{\color[HTML]{808080} \textit{Primario Esencial}}}                                                                                                                                                                                                                                                                                                             \\ \hline
        \multicolumn{1}{|l|}{Referencias}                            & \multicolumn{2}{l|}{}                                                                                                   & \multicolumn{5}{l|}{{\color[HTML]{808080} \textit{CU-15 Definir calendario}}}                                                                                                                                                                                                                                    \\ \hline
        \multicolumn{1}{|l|}{Precondición}                           & \multicolumn{7}{p{10cm}|}{{\color[HTML]{808080} \textit{}}}                                                                                                                                                                                                                                   \\ \hline
        \multicolumn{1}{|l|}{Poscondición}                           & \multicolumn{7}{p{10cm}|}{{\color[HTML]{808080} \textit{La actividad queda registrada en el sistema y disponible para completar su configuración (calendario, precios, etc.).}}}                                                                                                                                                                                                                                                       \\ \hline
        \multicolumn{1}{|l|}{Autor}                                  & \multicolumn{1}{l|}{{\color[HTML]{808080} \textit{Ismael Sallami Moreno}}}                           & \multicolumn{2}{l|}{Fecha}                                                   & \multicolumn{1}{l|}{}                          & \multicolumn{2}{l|}{Versión}                                                      & \multicolumn{1}{l|}{1.0}                                                     \\ \hline
        \end{tabular}
        \end{table}

        \begin{table}[H]
            \begin{tabular}{llllllll}
        \hline
        \multicolumn{8}{|l|}{\cellcolor[HTML]{F2F2F2}Propósito}                                                                                                                                                                                                                                                                                                                                                                                           \\ \hline
        \multicolumn{8}{|p{15cm}|}{{\color[HTML]{808080} \textit{Permitir al actor “Empleado” registrar una nueva actividad turística en el sistema, introduciendo todos los datos principales necesarios para su futura gestión y reserva por parte de los clientes.}}}                                                                                                                                                                                                             \\ \hline
        \end{tabular}
        \end{table}

        \begin{table}[H]
            \begin{tabular}{llllllll}
        \hline
        \multicolumn{8}{|l|}{\cellcolor[HTML]{F2F2F2}Resumen}                                                                                                                                                                                                                                                                                                                                                                                             \\ \hline
        \multicolumn{8}{|p{15cm}|}{{\color[HTML]{808080} \textit{Este caso de uso permite crear actividades turísticas que serán ofrecidas a los clientes en la aplicación. El empleado introduce los datos básicos de la actividad como el nombre, tipo (monumento, experiencia local, etc.), descripción, duración, localización y otros campos iniciales. Una vez creada, la actividad podrá ser completada con más detalles como calendarios, precios, condiciones de reserva, etc.}}} \\ \hline
        \end{tabular}
        \end{table}

    \begin{table}[H]
        \begin{tabular}{llllllll}
        \hline
        \multicolumn{8}{|l|}{\cellcolor[HTML]{F2F2F2}Curso Normal (Básico)}                                                                                                                                                                                                                                                                          \\ \hline
        \multicolumn{1}{|c|}{1}  & \multicolumn{1}{p{6cm}|}{{\color[HTML]{808080} \textit{El actor accede al formulario de creación de actividad, por ende, debe de acceder al sistema.}}} & \multicolumn{1}{l|}{{\color[HTML]{808080} }}                    & \multicolumn{5}{l|}{{\color[HTML]{808080} }}                                                                                            \\ \hline
        \multicolumn{1}{|c|}{}   & \multicolumn{1}{l|}{{\color[HTML]{808080} \textit{}}}                                                 & \multicolumn{1}{c|}{{\color[HTML]{808080} \textit{\textbf{2}}}} & \multicolumn{5}{p{6cm}|}{{\color[HTML]{808080} \textit{Introduce los datos principales (nombre, tipo, descripción, duración, ubicación, etc.).}}}                         \\ \hline
        \multicolumn{1}{|l|}{}   & \multicolumn{1}{l|}{{\color[HTML]{808080} }}                                                          & \multicolumn{1}{l|}{{\color[HTML]{808080} \textbf{3}}}          & \multicolumn{5}{p{6cm}|}{{\color[HTML]{808080} \textit{El sistema valida los datos introducidos.}}} \\ \hline
        \multicolumn{1}{|l|}{}   & \multicolumn{1}{l|}{{\color[HTML]{808080} }}                                                          & \multicolumn{1}{c|}{{\color[HTML]{808080} \textit{4}}}          & \multicolumn{5}{p{6cm}|}{{\color[HTML]{808080} \textit{El sistema almacena la nueva actividad en la base de datos.}}}                   \\ \hline
        \multicolumn{1}{|l|}{}   & \multicolumn{1}{l|}{{\color[HTML]{808080} }}                                                          & \multicolumn{1}{l|}{{\color[HTML]{808080} 5}}                   & \multicolumn{5}{p{6cm}|}{{\color[HTML]{808080} El sistema informa al actor de la correcta creación.}}                                  \\ \hline
        \end{tabular}
        \end{table}

        \begin{table}[H]
            \begin{tabular}{llllllll}
        \hline
        \multicolumn{8}{|l|}{\cellcolor[HTML]{F2F2F2}Cursos Alternos}                                                                                                                                                                                                                                                                                \\ \hline
        \multicolumn{1}{|c|}{3a} & \multicolumn{7}{p{15cm}|}{{\color[HTML]{808080} Si los datos son incompletos o erróneos, el sistema muestra errores y solicita correcciones.}}                                                                                                                                      \\ \hline
        \multicolumn{1}{|c|}{4a} & \multicolumn{7}{p{15cm}|}{{\color[HTML]{808080} Si ocurre un fallo en el almacenamiento, el sistema informa del error y sugiere reintentar.}}                                                                                                                                      \\ \hline
        \end{tabular}
        \end{table}

            % Please add the following required packages to your document preamble:
    % \usepackage[table,xcdraw]{xcolor}
    % Beamer presentation requires \usepackage{colortbl} instead of \usepackage[table,xcdraw]{xcolor}
    \begin{table}[H]
        \begin{tabular}{|l
        >{\columncolor[HTML]{FFFFFF}}l l
        >{\columncolor[HTML]{FFFFFF}}l 
        >{\columncolor[HTML]{FFFFFF}}l 
        >{\columncolor[HTML]{FFFFFF}}l 
        >{\columncolor[HTML]{FFFFFF}}l 
        >{\columncolor[HTML]{FFFFFF}}l |}
        \hline
        \multicolumn{8}{|l|}{\cellcolor[HTML]{F2F2F2}Otros datos}                                                                                                                                                                                                                                                                                                        \\ \hline
        \multicolumn{1}{|l|}{Frecuencia esperada} & \multicolumn{1}{p{5cm}|}{\cellcolor[HTML]{FFFFFF}{\color[HTML]{808080} }} & \multicolumn{1}{l|}{Rendimiento} & \multicolumn{5}{p{5cm}|}{\cellcolor[HTML]{FFFFFF}{\color[HTML]{808080} }} \\ \hline
        \multicolumn{1}{|l|}{Importancia}         & \multicolumn{1}{p{5cm}|}{\cellcolor[HTML]{FFFFFF}{\color[HTML]{808080} }}   & \multicolumn{1}{l|}{Urgencia}    & \multicolumn{5}{p{5cm}|}{\cellcolor[HTML]{FFFFFF}{\color[HTML]{808080} }}                          \\ \hline
        \multicolumn{1}{|l|}{Estado}              & \multicolumn{1}{p{5cm}|}{\cellcolor[HTML]{FFFFFF}{\color[HTML]{808080} }}                   & \multicolumn{1}{l|}{Estabilidad} & \multicolumn{5}{p{5cm}|}{\cellcolor[HTML]{FFFFFF}{\color[HTML]{808080} }}                                            \\ \hline
        \end{tabular}
    \end{table}

    \begin{table}[H]
        \begin{tabular}{|l|}
        \hline
        \rowcolor[HTML]{E6E6E6} 
        \multicolumn{1}{|c|}{\textbf{Comentarios}} \\ \hline
        \multicolumn{1}{|p{15cm}|}{{\color[HTML]{808080} \textit{}}} \\ \hline
        \end{tabular}
    \end{table}





    % ======Tabla Caso de uso Número: 15 =========================
    \newpage
    \begin{table}[H]
        \begin{tabular}{llllllll}
        \hline
        \multicolumn{1}{|l|}{Caso de Uso}                            & \multicolumn{5}{l|}{{\color[HTML]{808080} \textit{Definir calendario}}}                                                                                                                                                                & \multicolumn{2}{l|}{\cellcolor[HTML]{E6E6E6}{\color[HTML]{808080} \textit{CU-15}}}                                                   \\ \hline
        \multicolumn{1}{|l|}{Actores}                                & \multicolumn{7}{l|}{{\color[HTML]{808080} \textit{Empleado(I)}}}                                                                                                                                                                                                                                                                                               \\ \hline
        \multicolumn{1}{|l|}{Tipo}                                   & \multicolumn{7}{l|}{{\color[HTML]{808080} \textit{Primario Esencial}}}                                                                                                                                                                                                                                                                                                             \\ \hline
        \multicolumn{1}{|l|}{Referencias}                            & \multicolumn{2}{l|}{}                                                                                                   & \multicolumn{5}{l|}{{\color[HTML]{808080} \textit{}}}                                                                                                                                                                                                                                    \\ \hline
        \multicolumn{1}{|l|}{Precondición}                           & \multicolumn{7}{p{10cm}|}{{\color[HTML]{808080} \textit{La actividad a la cual se le elabora el calendario debe de existir}}}                                                                                                                                                                                                                                   \\ \hline
        \multicolumn{1}{|l|}{Poscondición}                           & \multicolumn{7}{p{10cm}|}{{\color[HTML]{808080} \textit{El calendario queda registrado en el sistema y asociado a la actividad correspondiente, permitiendo la gestión de reservas en las fechas definidas.}}}                                                                                                                                                                                                                                                       \\ \hline
        \multicolumn{1}{|l|}{Autor}                                  & \multicolumn{1}{l|}{{\color[HTML]{808080} \textit{Ismael Sallami Moreno}}}                           & \multicolumn{2}{l|}{Fecha}                                                   & \multicolumn{1}{l|}{}                          & \multicolumn{2}{l|}{Versión}                                                      & \multicolumn{1}{l|}{1.0}                                                     \\ \hline
        \end{tabular}
    \end{table}

    \begin{table}[H]
        \begin{tabular}{llllllll}
    \hline
    \multicolumn{8}{|l|}{\cellcolor[HTML]{F2F2F2}Propósito}                                                                                                                                                                                                                                                                                                                                                                                           \\ \hline
    \multicolumn{8}{|p{15cm}|}{{\color[HTML]{808080} \textit{Permitir al actor “Empleado” o definir un calendario de fechas y horarios disponibles para una actividad previamente registrada, facilitando la gestión de reservas por parte de los clientes.}}}                                                                                                                                                                                                             \\ \hline
    \end{tabular}
    \end{table}

    \begin{table}[H]
        \begin{tabular}{llllllll}
    \hline
    \multicolumn{8}{|l|}{\cellcolor[HTML]{F2F2F2}Resumen}                                                                                                                                                                                                                                                                                                                                                                                             \\ \hline
    \multicolumn{8}{|p{15cm}|}{{\color[HTML]{808080} \textit{Este caso de uso permite al encargado establecer las fechas y horarios en los que una actividad estará disponible para los clientes. El calendario puede incluir configuraciones como días específicos, horarios, capacidad máxima por sesión, y otras restricciones necesarias para la correcta gestión de la actividad.}}} \\ \hline
    \end{tabular}
    \end{table}

    \begin{table}[H]
        \begin{tabular}{llllllll}
    \hline
    \multicolumn{8}{|l|}{\cellcolor[HTML]{F2F2F2}Curso Normal (Básico)}                                                                                                                                                                                                                                                                          \\ \hline
    \multicolumn{1}{|c|}{1}  & \multicolumn{1}{l|}{{\color[HTML]{808080} \textit{El actor accede al formulario de definición de calendario.}}} & \multicolumn{1}{l|}{{\color[HTML]{808080} }}                    & \multicolumn{5}{l|}{{\color[HTML]{808080} }}                                                                                            \\ \hline
    \multicolumn{1}{|c|}{2}   & \multicolumn{1}{p{6cm}|}{{\color[HTML]{808080} \textit{Introduce las fechas, horarios y capacidad máxima para la actividad.}}}                                                 & \multicolumn{1}{c|}{} & \multicolumn{5}{p{6cm}|}{{\color[HTML]{808080} \textit{}}}                         \\ \hline
    \multicolumn{1}{|l|}{}   & \multicolumn{1}{l|}{{\color[HTML]{808080} }}                                                          & \multicolumn{1}{l|}{{\color[HTML]{808080} \textbf{3}}}          & \multicolumn{5}{p{6cm}|}{{\color[HTML]{808080} \textit{El sistema valida los datos introducidos.}}} \\ \hline
    \multicolumn{1}{|l|}{}   & \multicolumn{1}{l|}{{\color[HTML]{808080} }}                                                          & \multicolumn{1}{c|}{{\color[HTML]{808080} \textit{4}}}          & \multicolumn{5}{p{6cm}|}{{\color[HTML]{808080} \textit{El sistema almacena el calendario en la base de datos y lo asocia a la actividad.}}}                   \\ \hline
    \multicolumn{1}{|l|}{}   & \multicolumn{1}{l|}{{\color[HTML]{808080} }}                                                          & \multicolumn{1}{l|}{{\color[HTML]{808080} 5}}                   & \multicolumn{5}{p{6cm}|}{{\color[HTML]{808080} El sistema informa al actor de la correcta definición del calendario.}}                                  \\ \hline
    \end{tabular}
    \end{table}

    \begin{table}[H]
        \begin{tabular}{llllllll}
    \hline
    \multicolumn{8}{|l|}{\cellcolor[HTML]{F2F2F2}Cursos Alternos}                                                                                                                                                                                                                                                                                \\ \hline
    \multicolumn{1}{|c|}{3a} & \multicolumn{7}{p{15cm}|}{{\color[HTML]{808080} Si los datos son incompletos o erróneos, el sistema muestra errores y solicita correcciones.}}                                                                                                                                      \\ \hline
    \multicolumn{1}{|c|}{4a} & \multicolumn{7}{p{15cm}|}{{\color[HTML]{808080} Si ocurre un fallo en el almacenamiento, el sistema informa del error y sugiere reintentar.}}                                                                                                                                      \\ \hline
    \end{tabular}
    \end{table}

    \begin{table}[H]
        \begin{tabular}{|l
        >{\columncolor[HTML]{FFFFFF}}l l
        >{\columncolor[HTML]{FFFFFF}}l 
        >{\columncolor[HTML]{FFFFFF}}l 
        >{\columncolor[HTML]{FFFFFF}}l 
        >{\columncolor[HTML]{FFFFFF}}l 
        >{\columncolor[HTML]{FFFFFF}}l |}
        \hline
        \multicolumn{8}{|l|}{\cellcolor[HTML]{F2F2F2}Otros datos}                                                                                                                                                                                                                                                                                                        \\ \hline
        \multicolumn{1}{|l|}{Frecuencia esperada} & \multicolumn{1}{p{5cm}|}{\cellcolor[HTML]{FFFFFF}{\color[HTML]{808080} }} & \multicolumn{1}{l|}{Rendimiento} & \multicolumn{5}{p{5cm}|}{\cellcolor[HTML]{FFFFFF}{\color[HTML]{808080} }} \\ \hline
        \multicolumn{1}{|l|}{Importancia}         & \multicolumn{1}{p{5cm}|}{\cellcolor[HTML]{FFFFFF}{\color[HTML]{808080} }}   & \multicolumn{1}{l|}{Urgencia}    & \multicolumn{5}{p{5cm}|}{\cellcolor[HTML]{FFFFFF}{\color[HTML]{808080} }}                          \\ \hline
        \multicolumn{1}{|l|}{Estado}              & \multicolumn{1}{p{5cm}|}{\cellcolor[HTML]{FFFFFF}{\color[HTML]{808080} }}                   & \multicolumn{1}{l|}{Estabilidad} & \multicolumn{5}{p{5cm}|}{\cellcolor[HTML]{FFFFFF}{\color[HTML]{808080}}}                                            \\ \hline
        \end{tabular}
    \end{table}

    \begin{table}[H]
        \begin{tabular}{|l|}
        \hline
        \rowcolor[HTML]{E6E6E6} 
        \multicolumn{1}{|c|}{\textbf{Comentarios}} \\ \hline
        \multicolumn{1}{|p{15cm}|}{{\color[HTML]{808080} \textit{}}} \\ \hline
        \end{tabular}
    \end{table}



    % ======Tabla Caso de uso Número: 16 =========================
    \newpage
    \begin{table}[H]
        \begin{tabular}{llllllll}
        \hline
        \multicolumn{1}{|l|}{Caso de Uso}                            & \multicolumn{5}{l|}{{\color[HTML]{808080} \textit{Modificar actividad}}}                                                                                                                                                                & \multicolumn{2}{l|}{\cellcolor[HTML]{E6E6E6}{\color[HTML]{808080} \textit{CU-16}}}                                                   \\ \hline
        \multicolumn{1}{|l|}{Actores}                                & \multicolumn{7}{l|}{{\color[HTML]{808080} \textit{Empleado(I)}}}                                                                                                                                                                                                                                                                                               \\ \hline
        \multicolumn{1}{|l|}{Tipo}                                   & \multicolumn{7}{l|}{{\color[HTML]{808080} \textit{Primario Esencial}}}                                                                                                                                                                                                                                                                                                             \\ \hline
        \multicolumn{1}{|l|}{Referencias}                            & \multicolumn{2}{l|}{}                                                                                                   & \multicolumn{5}{l|}{{\color[HTML]{808080} \textit{CU-15 Definir calendario}}}                                                                                                                                                                                                                                    \\ \hline
        \multicolumn{1}{|l|}{Precondición}                           & \multicolumn{7}{p{10cm}|}{{\color[HTML]{808080} \textit{Debe de existir una actividad previamente.}}}                                                                                                                                                                                                                                   \\ \hline
        \multicolumn{1}{|l|}{Poscondición}                           & \multicolumn{7}{p{10cm}|}{{\color[HTML]{808080} \textit{La actividad queda actualizada con los nuevos datos introducidos.}}}                                                                                                                                                                                                                                                       \\ \hline
        \multicolumn{1}{|l|}{Autor}                                  & \multicolumn{1}{l|}{{\color[HTML]{808080} \textit{Ismael Sallami Moreno}}}                           & \multicolumn{2}{l|}{Fecha}                                                   & \multicolumn{1}{l|}{}                          & \multicolumn{2}{l|}{Versión}                                                      & \multicolumn{1}{l|}{1.0}                                                     \\ \hline
        \end{tabular}
    \end{table}

    \begin{table}[H]
        \begin{tabular}{llllllll}
        \hline
        \multicolumn{8}{|l|}{\cellcolor[HTML]{F2F2F2}Propósito}                                                                                                                                                                                                                                                                                                                                                                                           \\ \hline
        \multicolumn{8}{|p{15cm}|}{{\color[HTML]{808080} \textit{Permitir al actor “Empleado” actualizar la información de una actividad ya registrada, corrigiendo o ajustando los datos principales de la misma (nombre, descripción, duración, localización, tipo, etc.) según sea necesario.}}}                                                                                                                                                                                                             \\ \hline
        \end{tabular}
    \end{table}

    \begin{table}[H]
        \begin{tabular}{llllllll}
        \hline
        \multicolumn{8}{|l|}{\cellcolor[HTML]{F2F2F2}Resumen}                                                                                                                                                                                                                                                                                                                                                                                             \\ \hline
        \multicolumn{8}{|p{15cm}|}{{\color[HTML]{808080} \textit{Este caso de uso permite modificar las actividades existentes en el sistema, ya sea para corregir errores, actualizar detalles o mejorar la información ofrecida al cliente. Solo ciertos campos podrán modificarse, y en algunos casos pueden existir restricciones si la actividad ya tiene reservas vinculadas.}}} \\ \hline
        \end{tabular}
    \end{table}

    \begin{table}[H]
        \begin{tabular}{llllllll}
        \hline
        \multicolumn{8}{|l|}{\cellcolor[HTML]{F2F2F2}Curso Normal (Básico)}                                                                                                                                                                                                                                                                          \\ \hline
        \multicolumn{1}{|c|}{1}  & \multicolumn{1}{l|}{{\color[HTML]{808080} \textit{El actor accede al listado de actividades.}}} & \multicolumn{1}{l|}{{\color[HTML]{808080} }}                    & \multicolumn{5}{l|}{{\color[HTML]{808080} }}                                                                                            \\ \hline
        \multicolumn{1}{|c|}{2}   & \multicolumn{1}{p{6cm}|}{{\color[HTML]{808080} \textit{Selecciona la actividad a modificar.}}}                                                 & \multicolumn{1}{c|}{} & \multicolumn{5}{p{6cm}|}{{\color[HTML]{808080} \textit{}}}                         \\ \hline
        \multicolumn{1}{|l|}{}   & \multicolumn{1}{l|}{{\color[HTML]{808080} }}                                                          & \multicolumn{1}{l|}{{\color[HTML]{808080} \textbf{3}}}          & \multicolumn{5}{p{6cm}|}{{\color[HTML]{808080} \textit{El sistema muestra los datos actuales.}}} \\ \hline
        \multicolumn{1}{|c|}{4}  & \multicolumn{1}{p{6cm}|}{{\color[HTML]{808080} \textit{El actor edita los campos necesarios.}}} & \multicolumn{1}{l|}{{\color[HTML]{808080} }}                    & \multicolumn{5}{l|}{{\color[HTML]{808080} }}                                                                                            \\ \hline
        \multicolumn{1}{|l|}{}   & \multicolumn{1}{l|}{{\color[HTML]{808080} }}                                                          & \multicolumn{1}{l|}{{\color[HTML]{808080} 5}}                   & \multicolumn{5}{p{6cm}|}{{\color[HTML]{808080} El sistema valida los cambios.}}                                  \\ \hline
        \multicolumn{1}{|l|}{}   & \multicolumn{1}{l|}{{\color[HTML]{808080} }}                                                          & \multicolumn{1}{l|}{{\color[HTML]{808080} 6}}                   & \multicolumn{5}{p{6cm}|}{{\color[HTML]{808080} El sistema guarda la información actualizada.}}                                  \\ \hline
        \end{tabular}
    \end{table}

    \begin{table}[H]
        \begin{tabular}{llllllll}
        \hline
        \multicolumn{8}{|l|}{\cellcolor[HTML]{F2F2F2}Cursos Alternos}                                                                                                                                                                                                                                                                                \\ \hline
        \multicolumn{1}{|c|}{5a} & \multicolumn{7}{p{15cm}|}{{\color[HTML]{808080} Si los datos introducidos no son válidos, el sistema muestra errores y solicita corrección.}}                                                                                                                                      \\ \hline
        \multicolumn{1}{|c|}{6a} & \multicolumn{7}{p{15cm}|}{{\color[HTML]{808080} Si ocurre un fallo en el almacenamiento, el sistema informa del error y sugiere reintentar.}}                                                                                                                                      \\ \hline
        \end{tabular}
    \end{table}

    \begin{table}[H]
        \begin{tabular}{|l
        >{\columncolor[HTML]{FFFFFF}}l l
        >{\columncolor[HTML]{FFFFFF}}l 
        >{\columncolor[HTML]{FFFFFF}}l 
        >{\columncolor[HTML]{FFFFFF}}l 
        >{\columncolor[HTML]{FFFFFF}}l 
        >{\columncolor[HTML]{FFFFFF}}l |}
        \hline
        \multicolumn{8}{|l|}{\cellcolor[HTML]{F2F2F2}Otros datos}                                                                                                                                                                                                                                                                                                        \\ \hline
        \multicolumn{1}{|l|}{Frecuencia esperada} & \multicolumn{1}{p{5cm}|}{\cellcolor[HTML]{FFFFFF}{\color[HTML]{808080} }} & \multicolumn{1}{l|}{Rendimiento} & \multicolumn{5}{p{5cm}|}{\cellcolor[HTML]{FFFFFF}{\color[HTML]{808080} }} \\ \hline
        \multicolumn{1}{|l|}{Importancia}         & \multicolumn{1}{p{5cm}|}{\cellcolor[HTML]{FFFFFF}{\color[HTML]{808080} }}   & \multicolumn{1}{l|}{Urgencia}    & \multicolumn{5}{p{5cm}|}{\cellcolor[HTML]{FFFFFF}{\color[HTML]{808080} }}                          \\ \hline
        \multicolumn{1}{|l|}{Estado}              & \multicolumn{1}{p{5cm}|}{\cellcolor[HTML]{FFFFFF}{\color[HTML]{808080} }}                   & \multicolumn{1}{l|}{Estabilidad} & \multicolumn{5}{p{5cm}|}{\cellcolor[HTML]{FFFFFF}{\color[HTML]{808080} }}                                            \\ \hline
        \end{tabular}
    \end{table}

    \begin{table}[H]
        \begin{tabular}{|l|}
        \hline
        \rowcolor[HTML]{E6E6E6} 
        \multicolumn{1}{|c|}{\textbf{Comentarios}} \\ \hline
        \multicolumn{1}{|p{15cm}|}{{\color[HTML]{808080} \textit{}}} \\ \hline
        \end{tabular}
    \end{table}




    % ======Tabla Caso de uso Número: 17 =========================
    \newpage
    \begin{table}[H]
        \begin{tabular}{llllllll}
        \hline
        \multicolumn{1}{|l|}{Caso de Uso}                            & \multicolumn{5}{l|}{{\color[HTML]{808080} \textit{Eliminar actividad}}}                                                                                                                                                                & \multicolumn{2}{l|}{\cellcolor[HTML]{E6E6E6}{\color[HTML]{808080} \textit{CU-17}}}                                                   \\ \hline
        \multicolumn{1}{|l|}{Actores}                                & \multicolumn{7}{l|}{{\color[HTML]{808080} \textit{Empleado(I)}}}                                                                                                                                                                                                                                                                                               \\ \hline
        \multicolumn{1}{|l|}{Tipo}                                   & \multicolumn{7}{l|}{{\color[HTML]{808080} \textit{Primario Esencial}}}                                                                                                                                                                                                                                                                                                             \\ \hline
        \multicolumn{1}{|l|}{Referencias}                            & \multicolumn{2}{l|}{}                                                                                                   & \multicolumn{5}{l|}{}                                                                                                                                                                                                                                    \\ \hline
        \multicolumn{1}{|l|}{Precondición}                           & \multicolumn{7}{p{10cm}|}{{\color[HTML]{808080} \textit{La actividad debe de existir previamente en el sistema.}}}                                                                                                                                                                                                                                   \\ \hline
        \multicolumn{1}{|l|}{Poscondición}                           & \multicolumn{7}{p{10cm}|}{{\color[HTML]{808080} \textit{La actividad queda eliminada del sistema y no estará disponible para reservas futuras.}}}                                                                                                                                                                                                                                                       \\ \hline
        \multicolumn{1}{|l|}{Autor}                                  & \multicolumn{1}{l|}{{\color[HTML]{808080} \textit{Ismael Sallami Moreno}}}                           & \multicolumn{2}{l|}{Fecha}                                                   & \multicolumn{1}{l|}{}                          & \multicolumn{2}{l|}{Versión}                                                      & \multicolumn{1}{l|}{1.0}                                                     \\ \hline
        \end{tabular}
    \end{table}

    \begin{table}[H]
        \begin{tabular}{llllllll}
        \hline
        \multicolumn{8}{|l|}{\cellcolor[HTML]{F2F2F2}Propósito}                                                                                                                                                                                                                                                                                                                                                                                           \\ \hline
        \multicolumn{8}{|p{15cm}|}{{\color[HTML]{808080} \textit{Permitir al actor “Empleado” eliminar una actividad turística previamente registrada en el sistema, asegurando que no tenga reservas activas asociadas.}}}                                                                                                                                                                                                             \\ \hline
        \end{tabular}
    \end{table}

    \begin{table}[H]
        \begin{tabular}{llllllll}
        \hline
        \multicolumn{8}{|l|}{\cellcolor[HTML]{F2F2F2}Resumen}                                                                                                                                                                                                                                                                                                                                                                                             \\ \hline
        \multicolumn{8}{|p{15cm}|}{{\color[HTML]{808080} \textit{Este caso de uso permite eliminar actividades que ya no serán ofrecidas a los clientes. Antes de proceder, el sistema verifica que no existan reservas activas asociadas a la actividad. Si las condiciones se cumplen, la actividad se elimina de la base de datos y deja de estar disponible en el sistema.}}} \\ \hline
        \end{tabular}
    \end{table}

    \begin{table}[H]
        \begin{tabular}{llllllll}
        \hline
        \multicolumn{8}{|l|}{\cellcolor[HTML]{F2F2F2}Curso Normal (Básico)}                                                                                                                                                                                                                                                                          \\ \hline
        \multicolumn{1}{|c|}{1}  & \multicolumn{1}{l|}{{\color[HTML]{808080} \textit{El actor accede al listado de actividades.}}} & \multicolumn{1}{l|}{{\color[HTML]{808080} }}                    & \multicolumn{5}{l|}{{\color[HTML]{808080} }}                                                                                            \\ \hline
        \multicolumn{1}{|c|}{2}   & \multicolumn{1}{p{6cm}|}{{\color[HTML]{808080} \textit{Selecciona la actividad a eliminar.}}}                                                 & \multicolumn{1}{c|}{} & \multicolumn{5}{p{6cm}|}{{\color[HTML]{808080} \textit{}}}                         \\ \hline
        \multicolumn{1}{|l|}{}   & \multicolumn{1}{l|}{{\color[HTML]{808080} }}                                                          & \multicolumn{1}{l|}{{\color[HTML]{808080} \textbf{3}}}          & \multicolumn{5}{p{6cm}|}{{\color[HTML]{808080} \textit{El sistema verifica que no existan reservas activas asociadas.}}} \\ \hline
        \multicolumn{1}{|c|}{}  & \multicolumn{1}{p{6cm}|}{{\color[HTML]{808080} }} & \multicolumn{1}{l|}{{\color[HTML]{808080} \textbf{4}}}                    & \multicolumn{5}{p{6cm}|}{{\color[HTML]{808080} \textit{El sistema solicita confirmación para eliminar la actividad.}}}                                                                                            \\ \hline
        \multicolumn{1}{|l|}{}   & \multicolumn{1}{l|}{{\color[HTML]{808080} }}                                                          & \multicolumn{1}{l|}{{\color[HTML]{808080} 5}}                   & \multicolumn{5}{p{6cm}|}{{\color[HTML]{808080} El sistema elimina la actividad de la base de datos.}}                                  \\ \hline
        \multicolumn{1}{|l|}{}   & \multicolumn{1}{l|}{{\color[HTML]{808080} }}                                                          & \multicolumn{1}{l|}{{\color[HTML]{808080} 6}}                   & \multicolumn{5}{p{6cm}|}{{\color[HTML]{808080} El sistema informa al actor de la correcta eliminación de la actividad.}}                                  \\ \hline
        \end{tabular}
    \end{table}

    \begin{table}[H]
        \begin{tabular}{llllllll}
        \hline
        \multicolumn{8}{|l|}{\cellcolor[HTML]{F2F2F2}Cursos Alternos}                                                                                                                                                                                                                                                                                \\ \hline
        \multicolumn{1}{|c|}{3a} & \multicolumn{7}{p{15cm}|}{{\color[HTML]{808080} Si existen reservas activas asociadas, el sistema informa al actor que no se puede eliminar la actividad.}}                                                                                                                                      \\ \hline
        \multicolumn{1}{|c|}{5a} & \multicolumn{7}{p{15cm}|}{{\color[HTML]{808080} Si ocurre un fallo en el almacenamiento, el sistema informa del error y sugiere reintentar.}}                                                                                                                                      \\ \hline
        \end{tabular}
    \end{table}
    \begin{table}[H]
        \begin{tabular}{|l
        >{\columncolor[HTML]{FFFFFF}}l l
        >{\columncolor[HTML]{FFFFFF}}l 
        >{\columncolor[HTML]{FFFFFF}}l 
        >{\columncolor[HTML]{FFFFFF}}l 
        >{\columncolor[HTML]{FFFFFF}}l 
        >{\columncolor[HTML]{FFFFFF}}l |}
        \hline
        \multicolumn{8}{|l|}{\cellcolor[HTML]{F2F2F2}Otros datos}                                                                                                                                                                                                                                                                                                        \\ \hline
        \multicolumn{1}{|l|}{Frecuencia esperada} & \multicolumn{1}{p{5cm}|}{\cellcolor[HTML]{FFFFFF}} & \multicolumn{1}{l|}{Rendimiento} & \multicolumn{5}{p{5cm}|}{\cellcolor[HTML]{FFFFFF}} \\ \hline
        \multicolumn{1}{|l|}{Importancia}         & \multicolumn{1}{p{5cm}|}{\cellcolor[HTML]{FFFFFF}}   & \multicolumn{1}{l|}{Urgencia}    & \multicolumn{5}{p{5cm}|}{\cellcolor[HTML]{FFFFFF}}                          \\ \hline
        \multicolumn{1}{|l|}{Estado}              & \multicolumn{1}{p{5cm}|}{\cellcolor[HTML]{FFFFFF}}                   & \multicolumn{1}{l|}{Estabilidad} & \multicolumn{5}{p{5cm}|}{\cellcolor[HTML]{FFFFFF}}                                            \\ \hline
        \end{tabular}
    \end{table}

    \begin{table}[H]
        \begin{tabular}{|l|}
        \hline
        \rowcolor[HTML]{E6E6E6} 
        \multicolumn{1}{|c|}{\textbf{Comentarios}} \\ \hline
        \multicolumn{1}{|p{15cm}|}{{\color[HTML]{808080} \textit{Podríamos considerar que el encargado actúa, pero en función del diagrama lo exponemos de esta manera, así como el empleado es quien inicia la acción.}}} \\ \hline
        \end{tabular}
    \end{table}





    
